\section{Microeconomic Model} \label{sec:micromodel}

This paper is based on a microeconomic model of an educational production function, in which a student's academic performance is the result of a combination of different inputs throughout the learning process. The approach follows the classical literature on the economics of education, in which learning is treated as the product of individual, family, school, and social factors \parencite{funcaodeproducaoeducacional}.

Formally, the academic outcome of student $i$ in period $t$, denoted by $A_{it}$, can be described by the following educational production function:

\begin{equation}
A_{it} = f \left( B_{i}^{(t)}, P_{i}^{(t)}, S_{i}^{(t)}, I_i \right).
\end{equation}

In this expression, $B_{i}^{(t)}$ represents the student's family \textit{background}, which includes socioeconomic characteristics of the household, cultural capital, and educational support conditions. $P_{i}^{(t)}$ corresponds to the set of peers with whom the student interacts in the school environment, reflecting possible educational externalities arising from the composition of the class or school.

The term $S_{i}^{(t)}$ represents the school inputs available to the student. In this paper, these inputs are modeled explicitly as a function of three main dimensions of the educational environment:

\begin{equation} 
S_{i}^{(t)} = S(\tau,\pi,\sigma)_{i}^{(t)}, 
\end{equation}

where $\tau$ represents the time the student spends in school, $\pi$ corresponds to the personal and socio-emotional development provided by the school environment, and $\sigma$ represents the infrastructure and pedagogical resources available. Thus, school inputs are understood as a combination of these three factors, which jointly affect the learning process.

Finally, $I_i$ represents the student's innate abilities—that is, relatively persistent individual characteristics that influence their learning capacity.

It is assumed that all these factors positively impact the student's academic outcome, so that the partial derivatives of the production function are positive. Formally,

$$\frac{\partial A_{it}}{\partial B_{i}^{(t)}} > 0,$$
$$\quad$$
$$\frac{\partial A_{it}}{\partial P_{i}^{(t)}} > 0,$$
$$\quad$$
$$\frac{\partial A_{it}}{\partial S(\tau,\pi,\sigma)_{i}^{(t)}} > 0,$$
$$\quad$$
$$\frac{\partial A_{it}}{\partial I_{i}} > 0.$$

Family \textit{background} $B_{i}^{(t)}$ positively impacts performance, as a better family context is associated with better educational support conditions, which raises the student's academic outcome. Peers $P_{i}^{(t)}$ also exert a positive effect, since learning presents positive externalities: better peers tend to generate an environment more conducive to learning, raising individual performance.

School inputs $S(\tau,\pi,\sigma)_{i}^{(t)}$ positively affect the result and are composed of three dimensions. First, the time spent in school $\tau$ has a positive effect, i.e.,

$$\frac{\partial A_{it}}{\partial \tau} > 0,$$

because the more time dedicated to schooling, the greater the accumulation of knowledge tends to be. Second, personal development $\pi$ also presents a positive effect,

$$\frac{\partial A_{it}}{\partial \pi} > 0,$$

since greater emotional and social stability favors the learning process. Finally, school infrastructure $\sigma$ contributes positively to performance,

$$\frac{\partial A_{it}}{\partial \sigma} > 0,$$

as better physical and pedagogical resources expand the learning channels available to the student.

Additionally, innate abilities $I_{i}$ positively influence academic performance,

$$\frac{\partial A_{it}}{\partial I_{i}} > 0,$$

since higher levels of individual aptitude tend to raise a given student's learning capacity.

%-----------------------------------------------------------------------
\subsection{Channels of action of the PEI}
%-----------------------------------------------------------------------

In the context of the integral model, the main channels of impact on academic performance are peers and school factors. Regarding peers $P_{i}^{(t)}$, full-time schools may attract more engaged students with better prior performance, while simultaneously potentially pushing away more vulnerable students or those who need to balance study and work. From a theoretical point of view, this dynamic raises the expected average performance, i.e.,

$$\frac{\partial A_{it}}{\partial P_{i}^{(t)}} > 0.$$

However, from an empirical and public policy evaluation standpoint, this same dynamic represents exactly the compositional-change bias discussed in recent literature \parencite{favaro_composicao_pei}. This phenomenon requires econometric tools that reduce the risk of confusing the accounting effect generated by student turnover with learning gains, which underpins the methodological choices detailed in the next section.

Furthermore, school factors $S_{i}^{(t)}$, which can be decomposed into time spent $\tau$, personal development $\pi$, and infrastructure $\sigma$, also constitute relevant channels of direct impact of the program. Regarding the time spent in school, it is expected that

$$\frac{\partial A_{it}}{\partial \tau} > 0,$$

since greater exposure to the school environment tends to increase the accumulation of learning, although, in isolation, this channel may eventually present a limited effect if it is not accompanied by pedagogical changes, as seen in the experience of \textcite{AlmeidaEtAl2016MaisEducacao}.

With respect to personal development $\pi$, a positive effect on performance is expected,

$$\frac{\partial A_{it}}{\partial \pi} > 0,$$

as greater emotional, social, and behavioral stability, stimulated by PEI components such as the ``Life Project'' and tutoring, favors the continuity and quality of the learning process \textcite{araujo2020_pernambuco_impacto_enem}. 

Finally, the improvement in school infrastructure $\sigma$ expands pedagogical possibilities and available teaching methods, generating new forms of learning. Thus, it is expected that

$$\frac{\partial A_{it}}{\partial \sigma} > 0.$$

In summary, the integral model acts through two distinct channels on academic performance. The direct channel operates via the strengthening of school inputs (time spent $\tau$, personal development $\pi$, and infrastructure $\sigma$), representing the learning gain that the policy intends to generate. The indirect channel operates via the alteration of the composition of peers $P_i^{(t)}$: by imposing an extended school day and a more demanding pedagogical model, the PEI selectively alters who remains enrolled in the treated school, attracting more engaged students and pushing away the most vulnerable.

From an impact evaluation standpoint, this distinction is fundamental. The parameter of interest in this paper, the Average Treatment Effect on the Treated (ATT), approximates the direct channel that the policy seeks to produce: the learning gain associated with changes in school time and organization. The peer channel, however, represents a composition problem. The observed increase in average proficiency among PEI schools may partially reflect a change in the student profile, not only pedagogical improvement. Formally, when the composition of $P_i^{(t)}$ changes systematically between the pre- and post-treatment periods, the stationarity assumption is violated, and estimators that do not correct for this change may confuse the compositional effect with the effect attributed to the policy.
